% GDZ PPN236018647,native531;HKUST532.
\clearpage\SourceLeaf{531}\ArithmeticHead{502}{L}{TAGEBUCH MIT ERLÄUTERUNGEN.}
\NumberedItem{31}
\indent Numeros fractionum inaequalium quarum denominatores certum limitem non superant ad numerum fractionum omnium quarum num[eratores] aut denom[inatores] sint diversi infra eundem limitem in infinito ut \(6:\pi\pi\).
\par\noindent\hfill[1796] Sept. 6.

\medskip\begingroup\small
\indent Vergl. die Bemerkung zu der Nr. 14, oben S. 495.
\par\endgroup

\NumberedItem{32}
\indent Si \(\displaystyle\int\frac{dx}{\surd(1-x^3)}\) stat[uatur] \(\Pi:x=z\) et \(x=\Phi:z\), erit
\[
\begin{aligned}
\Phi:z={}&z-\frac18z^4+\frac1{112}z^7-\frac1{1792}z^{10}+\frac3{1792.52}z^{13}\\
&-\frac{3.185}{1792.52.14.15.16}z^{16}\ \ldotp\,\ldotp\,\ldotp
\end{aligned}
\]
\noindent\hfill[1796] Sept. 9.

\medskip\begingroup\small
\indent Vergl. den art. [5.] der \textit{Exercitationes mathematicae}, oben S. 140 und die zugehörige Anmerkung S. 144. Im Zähler des Koeffizienten von \(z^{16}\) muß es statt 185 heißen 165\footnote{Der Koeffizient von \(z^{16}\) lautet nämlich, vergl. S. 140 letzte Zeile
\[
-\frac{1.3.3.5.5.9.9.11.11.15}{1.2.3.4.5.6.7.8.9.10.11.12.13.14.15}\mathpunct{\raisebox{.5ex}{\normalfont .}}
\]
}. Das hier betrachtete Integral kommt auch in der oben S. 152 abgedruckten Leisteaufzeichnung vor. In der Werke VIII, S. 93 abgedruckten, aus dem Jahre 1800 stammenden Aufzeichnung untersucht \textsc{Gauss} die Umkehrung des Integrals \(\displaystyle\int\frac{dx}{\sqrt{1+x^3}}\). Vergl. auch die folgende Notiz Nr. 33.
\par\noindent\hfill\textsc{Schlesinger.}
\par\endgroup

\NumberedItem{33}
\indent Si
\[
\Phi:\int\frac{dx}{\surd(1-x^n)}=x
\]
\noindent erit:
\[
\Phi:z=z-\frac{1.z^n}{2.n+1}A+\frac{n-1.z^n}{4.2n+1}B-\frac{nn-n-1.[z^n]}{2.n+1.3n+1}C\ \ldotp\,\ldotp\,\ldotp
\]
\noindent\hfill[1796] Sept. 14.

\medskip\begingroup\small
\indent Hier bedeutet (vergl. die Bemerkung zum art. [5.] der \textit{Exercitationes mathematicae}, oben S. 144) \(A\) das erste Glied der Reihe, also \(z\), \(B\) das zweite, also \(+\frac{1.z^n}{2(n+1)}.z\), \(C\) das dritte, also \(\frac{(n-1)z^n}{4(2n+1)}.\frac{1.z^n}{2(n+1)}.z\) u.s.w. Diese Schreibweise wird auf \textsc{Newton} zurückgeführt. Für \(n=3\) ergeben sich die in der vorhergehenden Nummer betrachteten Ausdrücke.
\par\noindent\hfill\textsc{Schlesinger.}
\par\endgroup
